\chapter{Optimisation under Constraint}\label{ch:b1:constrained-optimisation}

A student in Brenmoor has sixty hours a week to split between paid work in
a café and study. Each hour of work pays \dal{15}. A grant of \dal{200} a
week comes with a rule: once the student earns more than \dal{100} a week,
the grant shrinks by half a dalen for each further dalen earned, until
nothing is left. How many hours should the student work? The question has
the shape of almost every problem in the microeconomics of this book: a
person or a firm wants as much as possible of something, and a constraint
--- hours, income, land, a budget --- stops them from having it all. This
chapter gives the method that solves such problems, and the number it
produces along the way: the value of relaxing the constraint by one unit.

\begin{openingchart}{Hours worked in a year by a person in work, against labour
compensation per hour, 111 countries, 2019. Compensation is in 2021 dollars at
purchasing-power parity, on a logarithmic axis. Venezuela, at 5 cents an hour
in these data, is left out.}{Penn World Table 11.0; average annual hours worked
by persons engaged; compensation per hour = labour share $\times$ output-side real
GDP at current PPPs $\div$ (persons engaged $\times$ hours)}
\begin{tikzpicture}
\begin{semilogxaxis}[omincomex, xlabel={Compensation per hour (dollars, logarithmic axis)},
  ylabel={Hours worked per year}, xmin=0.5, xmax=100, ymin=1200, ymax=2600,
  height=0.55\linewidth]
  \addplot table[x=wage, y=hours] {figdata/bachelor-1/02-constrained-optimisation/out/hours_wage.csv};
\end{semilogxaxis}
\end{tikzpicture}
\end{openingchart}

Where an hour of work is paid more, people work fewer hours. The
correlation between hours and the logarithm of pay per hour across the 111
countries is $-0.53$, and the fitted line falls by 162 hours a year for
each unit of the logarithm, about 1.6 hours a year less for each per cent
more pay. % figdata: hours_wage_stats
In the tenth of countries with the lowest pay (a median of \usd{2.20} an
hour) the median worker puts in 2064 hours a year; in the tenth with the
highest (\usd{49.48}) the median is 1504 hours. % figdata: hours_wage_tenths
In the United States, at \usd{45.86} an hour, it is 1797. % figdata: hours_wage_stats
A higher wage makes an hour of leisure more expensive, which should make
people work more, and it makes them richer, which should make them take
more leisure. The chart says that across countries the second force wins,
or that something else differs between rich and poor countries. To say which,
one must first know how a person chooses hours under a constraint. That is
the subject of this chapter, and \cref{ch:b1:labour-supply-and-saving}
returns to the chart with the answer.

\begin{recall}
For a function $f(x, y)$, the partial derivative $f_x = \pd{f}{x}$ is the
derivative in $x$ with $y$ held fixed. The gradient $\nabla f = (f_x, f_y)$
points in the direction of steepest increase. A level curve is the set of
points where $f$ takes one value. A critical point, where $\nabla f = 0$, is
a candidate for an unconstrained maximum.
\end{recall}

\section{Functions of two variables}\label{sec:b1:constrained-optimisation:two-variables}

An objective in economics usually depends on two or more quantities: the
amounts of two goods, of two inputs, of work and leisure. Its level curves
--- indifference curves for a consumer, isoquants for a firm --- carry most
of the reasoning, and the slope of a level curve is the number used most.

\begin{proposition}[Slope of a level curve]\label{prop:b1:constrained-optimisation:level-slope}
\index{level curve!slope}
Let $f$ have continuous partial derivatives, and let the level curve
$f(x, y) = c$ be, near a point where $f_y \neq 0$, the graph of a
differentiable function $y(x)$. Its slope there is
\[
y'(x) = -\frac{f_x}{f_y}.
\]
\end{proposition}

\begin{proof}
Along the curve, $f(x, y(x)) = c$ for every $x$. Differentiate with respect
to $x$ by the chain rule: $f_x + f_y\, y'(x) = 0$. Divide by $f_y \neq 0$.
\end{proof}

That the curve is the graph of such a function whenever $f_y \neq 0$ is the
implicit function theorem of the mathematics volumes; it is admitted here.
The proposition says that moving along a level curve, a small increase
$\dd x$ must be paid for with $\dd y = -(f_x/f_y)\dd x$: the ratio of the
partial derivatives is the rate at which the function lets one variable be
traded for the other. For $f(x, y) = x^{1/2}y^{1/2}$,
$f_x = \tfrac12 x^{-1/2}y^{1/2}$ and $f_y = \tfrac12 x^{1/2}y^{-1/2}$, so the
slope of a level curve is $-y/x$. The gradient $(f_x, f_y)$ is
perpendicular to the level curve, since it is perpendicular to the tangent
vector $(1, y')$: $f_x + f_y y' = 0$.

\section{A maximum under one equality constraint}\label{sec:b1:constrained-optimisation:equality}

\begin{definition}[Constrained maximum and minimum]\label{def:b1:constrained-optimisation:constrained-max}
\index{constrained maximum}\index{constrained minimum}
Let $f$ and $g$ be functions of $(x, y)$ and $b$ a number. A point
$(x^*, y^*)$ with $g(x^*, y^*) = b$ is a \emph{constrained maximum} of $f$
under the constraint $g = b$ if $f(x^*, y^*) \ge f(x, y)$ for every $(x, y)$
with $g(x, y) = b$ (in the region where the problem is posed). A
\emph{constrained minimum} reverses the inequality. The problem is written
$\max_{x, y} f(x, y)$ subject to $g(x, y) = b$.
\end{definition}

The first method is substitution: solve the constraint for one variable and
maximise a function of the other, with the calculus of one variable.

\begin{example}[Substitution]\label{ex:b1:constrained-optimisation:substitution}
Maximise $f(x, y) = x^{1/2}y^{1/2}$ subject to $x + 2y = 12$, with $x, y > 0$.
The constraint gives $x = 12 - 2y$. Since the square root is increasing,
maximising $f$ is maximising $h(y) = (12 - 2y)y = 12y - 2y^2$:
$h'(y) = 12 - 4y = 0$ at $y = 3$, and $h'' = -4 < 0$, so this is the
maximum. Then $x = 6$ and $f^* = \sqrt{18} = 4.243$.
\end{example}

Substitution works when the constraint can be solved, and it hides the
geometry. The geometry is this: at the best point on the constraint, the
level curve of $f$ through that point cannot cross the constraint, for if
it did, a step along the constraint would reach a higher level curve.

\begin{definition}[Tangency condition]\label{def:b1:constrained-optimisation:tangency}
\index{tangency condition}
For the problem of maximising $f$ subject to $g = b$, the \emph{tangency
condition} at a point is that the level curve of $f$ and the constraint
curve have the same slope there:
\[
\frac{f_x}{f_y} = \frac{g_x}{g_y}.
\]
\end{definition}

\begin{proposition}[Tangency at an interior optimum]\label{prop:b1:constrained-optimisation:tangency}
\index{tangency condition!necessity}
Let $f$ and $g$ have continuous partial derivatives, let $(x^*, y^*)$ be a
\omterm{def:b1:constrained-optimisation:constrained-max}{constrained maximum} of $f$ subject to $g = b$ that is not at the edge of the
region, and let $g_y \neq 0$ and $f_y \neq 0$ there. Then the \omterm{def:b1:constrained-optimisation:tangency}{tangency
condition} holds at $(x^*, y^*)$.
\end{proposition}

\begin{proof}
Near $(x^*, y^*)$ the constraint is the graph of a function $y(x)$ with
$y' = -g_x/g_y$ (\cref{prop:b1:constrained-optimisation:level-slope} applied
to $g$). The function $h(x) = f(x, y(x))$ has a maximum at $x^*$, an interior
point, so $h'(x^*) = f_x + f_y y' = f_x - f_y g_x/g_y = 0$. Divide by
$f_y \neq 0$: $f_x/f_y = g_x/g_y$.
\end{proof}

In \cref{ex:b1:constrained-optimisation:substitution}, the slope of the
level curve at $(6, 3)$ is $-y/x = -1/2$ and the slope of the constraint
$x + 2y = 12$ is $-1/2$: tangency.

\begin{omfigure}
\begin{tikzpicture}
\begin{axis}[omrecord, grid=none, xlabel={$x$}, ylabel={$y$}, xmin=0, xmax=13,
  ymin=0, ymax=7, width=0.62\linewidth, height=0.44\linewidth, samples=120,
  clip mode=individual]
  \addplot[demand, domain=1.3:13] {9/x};
  \addplot[demand, domain=2.6:13] {18/x};
  \addplot[demand, domain=3.6:13] {25/x};
  \addplot[supply, domain=0:12] {(12 - x)/2};
  \addplot[guide] coordinates {(6,0) (6,3)};
  \addplot[guide] coordinates {(0,3) (6,3)};
  \addplot[only marks, mark=*, black] coordinates {(6,3)};
  \node[font=\footnotesize, anchor=west] at (axis cs:13,1.385) {$f = \sqrt{18}$};
  \node[font=\footnotesize, anchor=west] at (axis cs:13,1.923) {$f = 5$};
  \node[font=\footnotesize, anchor=west] at (axis cs:13,0.692) {$f = 3$};
  \node[font=\footnotesize, anchor=west, text=omThm] at (axis cs:4.3,6.4) {constraint $x + 2y = 12$};
\end{axis}
\end{tikzpicture}

{\small Level curves of $f = x^{1/2}y^{1/2}$ (blue) and the constraint
$x + 2y = 12$ (red). The highest level curve that the constraint reaches,
$f = \sqrt{18}$, touches it at $(6, 3)$; the curve $f = 5$ lies beyond reach.}
\end{omfigure}

\section{Lagrange's method}\label{sec:b1:constrained-optimisation:lagrange}

The \omterm{def:b1:constrained-optimisation:tangency}{tangency condition} and the constraint are two equations in two
unknowns. \omterm{met:b1:constrained-optimisation:lagrange}{Lagrange's method} writes them in a form that extends to many
variables and many constraints, and that delivers one more number.

\begin{definition}[Lagrangian]\label{def:b1:constrained-optimisation:lagrangian}
\index{Lagrangian}
For the problem $\max f(x, y)$ subject to $g(x, y) = b$, the
\emph{Lagrangian} is the function of three variables
\[
\mathcal{L}(x, y, \lambda) = f(x, y) - \lambda\bigl(g(x, y) - b\bigr),
\]
where the new variable $\lambda$ is the Lagrange multiplier.
\end{definition}

\begin{theorem}[Lagrange conditions]\label{thm:b1:constrained-optimisation:lagrange}
\index{Lagrange conditions}
Let $f$ and $g$ have continuous partial derivatives on an open set, and let
$(x^*, y^*)$ be a \omterm{def:b1:constrained-optimisation:constrained-max}{constrained maximum} (or minimum) of $f$ subject to $g = b$
at which $\nabla g \neq 0$. Then there is a number $\lambda^*$ such that
$(x^*, y^*, \lambda^*)$ is a critical point of $\mathcal L$:
\[
f_x = \lambda g_x, \qquad f_y = \lambda g_y, \qquad g(x, y) = b .
\]
The same holds with $n$ variables and $m < n$ constraints, one multiplier
per constraint, when the gradients of the constraints are linearly
independent.
\end{theorem}
\admitted

\begin{proposition}[Lagrange and tangency]\label{prop:b1:constrained-optimisation:equivalence}
\index{Lagrange conditions!and tangency}
At a point where $f_y \neq 0$ and $g_y \neq 0$, the \omterm{thm:b1:constrained-optimisation:lagrange}{Lagrange conditions} hold
for some $\lambda$ if and only if the constraint and the \omterm{def:b1:constrained-optimisation:tangency}{tangency condition}
hold; then $\lambda = f_y/g_y = f_x/g_x$ (the last when $g_x \neq 0$).
\end{proposition}

\begin{proof}
If $f_x = \lambda g_x$ and $f_y = \lambda g_y$, then $\lambda = f_y/g_y \neq 0$
and $f_x/f_y = \lambda g_x/(\lambda g_y) = g_x/g_y$. Conversely, if
$f_x/f_y = g_x/g_y$, set $\lambda = f_y/g_y$: the second condition holds by
construction, and $f_x = f_y g_x/g_y = \lambda g_x$ gives the first.
\end{proof}

\begin{method}[Solving a problem with Lagrange's method]\label{met:b1:constrained-optimisation:lagrange}
\index{Lagrange's method}
\begin{enumerate}
  \item Write the \omterm{def:b1:constrained-optimisation:lagrangian}{Lagrangian} $\mathcal L = f - \lambda(g - b)$.
  \item Set its three partial derivatives to zero.
  \item Eliminate $\lambda$ by dividing the first two conditions: this gives
    the \omterm{def:b1:constrained-optimisation:tangency}{tangency condition}, a relation between $x$ and $y$.
  \item Substitute that relation into the constraint, solve, and recover
    $\lambda$ from either of the first two conditions.
  \item Check that the point is a maximum: compare its value with the edges
    of the region (\cref{sec:b1:constrained-optimisation:corners}).
\end{enumerate}
\end{method}

For \cref{ex:b1:constrained-optimisation:substitution}: $\tfrac12\sqrt{y/x}
= \lambda$ and $\tfrac12\sqrt{x/y} = 2\lambda$. Dividing the second by the
first, $x/y = 2$, so $x = 2y$; the constraint gives $4y = 12$, $y = 3$,
$x = 6$, and $\lambda = \tfrac12\sqrt{3/6} = 0.354$. The same point, found
without solving the constraint, and a new number $\lambda$.

\section{The multiplier as a shadow price}\label{sec:b1:constrained-optimisation:shadow}

\begin{definition}[Value function]\label{def:b1:constrained-optimisation:value-function}
\index{value function}
The \emph{value function} of the problem $\max f$ subject to $g = b$ is the
best attainable value as a function of the constraint's level,
$V(b) = f\bigl(x^*(b), y^*(b)\bigr)$, where $(x^*(b), y^*(b))$ is the solution
for that $b$.
\end{definition}

\begin{proposition}[The multiplier is the derivative of the value function]\label{prop:b1:constrained-optimisation:shadow}
\index{shadow price!as a derivative}
If the solution $(x^*(b), y^*(b))$ and its multiplier $\lambda^*(b)$ satisfy
the \omterm{thm:b1:constrained-optimisation:lagrange}{Lagrange conditions} and $x^*$, $y^*$ are differentiable in $b$, then
\[
V'(b) = \lambda^*(b).
\]
\end{proposition}

\begin{proof}
Differentiate $V(b) = f(x^*(b), y^*(b))$:
$V'(b) = f_x \,x^{*\prime} + f_y \,y^{*\prime}$. By the \omterm{thm:b1:constrained-optimisation:lagrange}{Lagrange conditions}
this is $\lambda^*(g_x x^{*\prime} + g_y y^{*\prime})$. The bracket is the
derivative of $g(x^*(b), y^*(b))$, which equals $b$ for every $b$, so the
bracket is $1$.
\end{proof}

\begin{definition}[Shadow price]\label{def:b1:constrained-optimisation:shadow-price}
\index{shadow price}\index{Lagrange multiplier!as a shadow price}
The \emph{shadow price} of a constraint is the value of its Lagrange
multiplier at the optimum, $\lambda^* = V'(b)$: the increase in the best
attainable value of the objective per unit by which the constraint is
relaxed. Its unit is the unit of the objective divided by the unit of the
constraint.
\end{definition}

In the example, $x^* = b/2$ and $y^* = b/4$, so
$V(b) = \sqrt{b^2/8} = b/(2\sqrt2)$ and $V'(b) = 1/(2\sqrt 2) = 0.354$, the
$\lambda$ found above. Relaxing the constraint from 12 to 14 raises the
value by $2 \times 0.354 = 0.71$, from $4.243$ to $4.950$; here the
approximation is exact because $V$ is linear in $b$. In general $\lambda$
gives the first-order change, $\Delta V \approx \lambda \Delta b$, as any
derivative does.

\begin{omfigure}
\begin{tikzpicture}
\begin{axis}[omrecord, grid=none, xlabel={$x$}, ylabel={$y$}, xmin=0, xmax=15,
  ymin=0, ymax=8, width=0.62\linewidth, height=0.44\linewidth, samples=120]
  \addplot[demand, domain=2.6:15] {18/x};
  \addplot[demand, shifted, domain=3.2:15] {24.5/x};
  \addplot[supply, domain=0:12] {(12 - x)/2};
  \addplot[supply, shifted, domain=0:14] {(14 - x)/2};
  \addplot[only marks, mark=*, black] coordinates {(6,3) (7,3.5)};
  \node[font=\footnotesize, anchor=north east] at (axis cs:6,3) {$V(12) = 4.243$};
  \node[font=\footnotesize, anchor=south west] at (axis cs:7.2,3.6) {$V(14) = 4.950$};
\end{axis}
\end{tikzpicture}

{\small The \omterm{def:b1:constrained-optimisation:shadow-price}{shadow price}: relaxing the constraint from $b = 12$ to $b = 14$
(dashed) moves the optimum from $(6, 3)$ to $(7, 3.5)$ and raises the value by
$\lambda^* \Delta b = 0.354 \times 2 = 0.71$.}
\end{omfigure}

The unit makes the multiplier economics. If $f$ is a firm's profit in
dalens and $b$ is hectares of land, $\lambda^*$ is in dalens per hectare: the
most the firm would pay to rent one more hectare, whatever the market rent
happens to be. If $f$ is utility and $b$ is income, $\lambda^*$ is utility
per dalen, the marginal utility of income. If the constraint is a
regulation --- a cap on emissions, a quota --- its \omterm{def:b1:constrained-optimisation:shadow-price}{shadow price} is what the
regulation costs the one who obeys it, at the margin.
\Cref{ch:b1:comparative-statics} extends the proposition to the derivative
of a \omterm{def:b1:constrained-optimisation:value-function}{value function} with respect to any parameter.

\section{Corners and kinks}\label{sec:b1:constrained-optimisation:corners}

Quantities in economics cannot be negative, and many constraints are
inequalities: a budget need not be spent, a field need not be fully
planted.

\begin{definition}[Binding constraint]\label{def:b1:constrained-optimisation:binding}
\index{binding constraint}
An inequality constraint $g(x, y) \le b$ is \emph{binding} at a solution if
it holds with equality there, $g(x^*, y^*) = b$, and \emph{slack} if it
holds strictly. A slack constraint has a \omterm{def:b1:constrained-optimisation:shadow-price}{shadow price} of zero: relaxing it a
little changes nothing.
\end{definition}

\begin{definition}[Interior and corner solutions]\label{def:b1:constrained-optimisation:corner}
\index{corner solution}\index{interior solution}
A solution with every quantity strictly positive is an \emph{interior
solution}. A solution at which some quantity is zero because a
non-negativity constraint binds is a \emph{corner solution}.
\end{definition}

At a corner the \omterm{def:b1:constrained-optimisation:tangency}{tangency condition} need not hold, and usually does not.
Take $f(x, y) = x + \ln y$ with the constraint $x + y = m$ and $x, y \ge 0$.
\omterm{met:b1:constrained-optimisation:lagrange}{Lagrange's method} gives $1 = \lambda$ and $1/y = \lambda$, so $y = 1$ and
$x = m - 1$. When $m < 1$ this asks for a negative $x$, which is not
allowed. The solution is the corner $x = 0$, $y = m$: every dalen goes to $y$,
because the last unit of $y$ is worth $1/m > 1$, more than a unit of $x$.
The rule at a corner is an inequality: the good that is not bought is worth
less, at the margin, than its price in terms of the good that is bought.

\begin{omfigure}
\begin{tikzpicture}
\begin{axis}[omrecord, grid=none, xlabel={$x$}, ylabel={$y$}, xmin=0, xmax=1.6,
  ymin=0, ymax=1.6, width=0.5\linewidth, height=0.42\linewidth, samples=100]
  \addplot[demand, domain=0:1.6] {exp(ln(0.5) - x)};
  \addplot[demand, domain=0:1.6] {exp(ln(1.0) - x)};
  \addplot[demand, domain=0:1.6] {exp(ln(1.5) - x)};
  \addplot[supply, domain=0:0.5] {0.5 - x};
  \addplot[only marks, mark=*, black] coordinates {(0,0.5)};
  \node[font=\footnotesize, anchor=west] at (axis cs:0.08,0.7) {corner $(0, 0.5)$};
\end{axis}
\end{tikzpicture}

{\small A corner: the level curves of $x + \ln y$ (blue) and the constraint
$x + y = 0.5$ (red). Every point of the constraint with $x > 0$ lies on a
lower level curve than the corner $(0, 0.5)$, where the level curve is flatter
than the constraint (slope $-0.5$ against $-1$): a unit of $x$ is worth only half
a unit of $y$ there.}
\end{omfigure}

The first-order conditions find candidates; they do not certify a maximum.
A level curve can touch the constraint from the wrong side, as the level
curves of $x^2 + y^2$ touch a straight line: there the tangency point is the
\omterm{def:b1:constrained-optimisation:constrained-max}{constrained minimum}. A sufficient condition with two variables is that the
level curves of $f$ are convex towards the origin --- the upper level sets
$\{f \ge c\}$ are convex --- and the constraint is a straight line; the
general test, with a bordered matrix of second derivatives, is in the Year~2
mathematics volume. In this book, objectives will be chosen so that their
upper level sets are convex, and the edges are checked by hand.

Return to the student. With $h$ hours of work, weekly income is
$y = 15h + \max\bigl(0, 200 - 0.5\max(0, 15h - 100)\bigr)$. Below
$h = 100/15 = 6.67$ hours each hour adds \dal{15}. From there until the grant
is gone, at $15h - 100 = 400$, that is $h = 33.33$, each hour adds
$15 \times 0.5 = \num{7.5}$~dalens. After that, \dal{15} again. Income at the two
kinks is \dal{300} and \dal{500}. Drawn against hours of study, $60 - h$, the
constraint is a line with two kinks.

\begin{omfigure}
\begin{tikzpicture}
\begin{axis}[omrecord, grid=none, xlabel={Hours of study a week}, ylabel={Income (dalens a week)},
  xmin=0, xmax=62, ymin=0, ymax=950, width=0.68\linewidth, height=0.44\linewidth]
  \addplot[supply, domain=60-20/3:60] {200 + 15*(60 - x)};
  \addplot[supply, domain=60-100/3:60-20/3] {250 + 7.5*(60 - x)};
  \addplot[supply, domain=0:60-100/3] {15*(60 - x)};
  \addplot[guide, domain=0:60] {900 - 15*x};
  \addplot[only marks, mark=*, black] coordinates {(53.333,300) (26.667,500)};
  \node[font=\footnotesize, anchor=south west] at (axis cs:53.3,300) {300};
  \node[font=\footnotesize, anchor=south west] at (axis cs:26.7,510) {500};
\end{axis}
\end{tikzpicture}

{\small The student's weekly constraint (red): \dal{15} an hour, a grant of
\dal{200} withdrawn at half a dalen per dalen earned above \dal{100}. On the
middle segment an hour of work pays \dal{7.5}. The dotted line is the
constraint with no grant.}
\end{omfigure}

A kink where the reward for one more hour falls has a sharp prediction.
At the first kink, where the grant starts to be withdrawn, pay per hour
drops from \dal{15} to \dal{7.5}. Every student who would work past the kink
at \dal{15} an hour, but would not reach it at \dal{7.5}, stops exactly there:
students with different tastes end up at the same point. The model
predicts \emph{bunching}: more people exactly at such a kink than at nearby
points. At the second kink, where the grant is exhausted and pay per hour
rises back to \dal{15}, the prediction is the opposite: nobody stops there.

\begin{model}{The constrained optimiser}
\assumes An agent knows its objective and its constraint, and chooses the
point of the constraint with the highest value of the objective.
\predicts At an interior optimum the agent's rate of trade-off equals the
constraint's; a slack constraint has a \omterm{def:b1:constrained-optimisation:shadow-price}{shadow price} of zero; a kink where
the marginal reward falls attracts a bunch of agents with different tastes.
\tested The fall of hours with pay across countries (correlation $-0.53$
with log pay, 111 countries) is consistent with it but not a test: the model
predicts no sign there. % figdata: hours_wage_stats
The test is bunching: American tax filers with self-employment income bunch at the first kink of
the earned income tax credit (see the record box). % ledger: b1:saez2010-firstkink
\breaks When agents do not know the constraint: wage earners show no bunching
at the same kink, and the response to the credit is larger where knowledge of
its schedule is higher (Chetty, Friedman and Saez 2013); % ledger: b1:saez2010-firstkink, b1:cfs2013-knowledge
and when hours are set by the employer, not chosen.
\end{model}

\begin{record}[Bunching at a kink]
The American earned income tax credit pays low earners a credit that rises
with earnings, then stays flat, then is withdrawn: its schedule has kinks.
Pooling the tax returns of 1995 to 2004, Saez (2010) found clear bunching of
reported earnings at the first kink, where the credit stops rising. The
implied elasticity of reported earnings was 0.21 for filers with one child
and 0.15 for those with two or more. % ledger: b1:saez2010-firstkink
The whole effect came from the self-employed, with elasticities of 1.1 and
0.8; filers with wage income only showed elasticities close to zero. % ledger: b1:saez2010-firstkink
At the threshold where income tax begins, bunching was clearest in the
1960s, with an elasticity around 0.2, and there was none at the higher
kinks of the schedule. % ledger: b1:saez2010-bracket
The prediction holds where the agent can adjust the number reported, and
fails where hours are hard to adjust or the schedule is not known. The self-employed
report their own earnings; a wage earner's are reported by the employer.
\end{record}

\section*{Exercises}\label{sec:b1:constrained-optimisation:exercises}

\begin{exercise}[$\star$]\label{exo:b1:constrained-optimisation:read-chart}
From the table behind the opening chart: in the tenth of countries with the
lowest pay, the median pay is \usd{2.20} an hour and the median hours 2064;
in the top tenth, \usd{49.48} and 1504. % figdata: hours_wage_tenths
\begin{enumerate}
  \item By what factor does pay differ, and by how many hours a year, and by what percentage, do hours differ?
  \item How many fewer hours does that make per week, over 48 working weeks?
  \item Is this the slope of one worker's labour supply? Say why not in one sentence.
\end{enumerate}
\end{exercise}

\begin{exercise}[$\star$]\label{exo:b1:constrained-optimisation:substitution}
Maximise $xy$ subject to $2x + 3y = 24$, $x, y > 0$, by substitution. Check
the \omterm{def:b1:constrained-optimisation:tangency}{tangency condition} at your answer.
\end{exercise}

\begin{exercise}[$\star$]\label{exo:b1:constrained-optimisation:level-slope}
Compute the slope of the level curve of $f(x, y) = x^2 y$ through the point
$(1, 2)$. Which level is it, and what is the equation of the curve?
\end{exercise}

\begin{exercise}[$\star$]\label{exo:b1:constrained-optimisation:binding}
The solution of $\max x^{1/2}y^{1/2}$ subject to $x + 2y \le 12$, $x \le 8$ and
$y \le 5$ is $(6, 3)$. Which constraints bind, and which have a \omterm{def:b1:constrained-optimisation:shadow-price}{shadow price}
of zero? What happens to the solution if the second constraint becomes
$x \le 4$?
\end{exercise}

\begin{exercise}[$\star\star$]\label{exo:b1:constrained-optimisation:lagrange-log}
Maximise $\ln x + 2\ln y$ subject to $p_1 x + p_2 y = m$ with \omterm{met:b1:constrained-optimisation:lagrange}{Lagrange's
method}. Give $x^*$, $y^*$ and $\lambda^*$ as functions of $p_1$, $p_2$, $m$, then
evaluate them at $p_1 = 1$, $p_2 = 2$, $m = 30$.
\end{exercise}

\begin{exercise}[$\star\star$]\label{exo:b1:constrained-optimisation:lambda-units}
In \cref{exo:b1:constrained-optimisation:lagrange-log}, at $p_1 = 1$,
$p_2 = 2$, $m = 30$:
\begin{enumerate}
  \item give the unit of $\lambda^*$ and say in words what it measures;
  \item predict, with $\lambda^*$, the change in the maximum value if $m$ rises to 40;
  \item compute the exact change and explain the gap.
\end{enumerate}
\end{exercise}

\begin{exercise}[$\star\star$]\label{exo:b1:constrained-optimisation:corner}
Maximise $x + \ln y$ subject to $x + y = 0.5$, $x, y \ge 0$. Show that
\omterm{met:b1:constrained-optimisation:lagrange}{Lagrange's method} gives an inadmissible point, find the solution, and
compute the ratio of the marginal values of $y$ and $x$ there.
\end{exercise}

\begin{exercise}[$\star\star$]\label{exo:b1:constrained-optimisation:student-kinks}
For the student of the chapter:
\begin{enumerate}
  \item compute the hours at which the two kinks occur and the income there;
  \item compute what an hour of work adds to income on each of the three segments;
  \item compute weekly income at 20 hours and at 40 hours of work.
\end{enumerate}
\end{exercise}

\begin{exercise}[$\star\star$]\label{exo:b1:constrained-optimisation:withdrawal}
The grant's withdrawal rate rises from 0.5 to 0.75, the other rules
unchanged. Predict, with reasons, the direction of the change in hours for
three students: one who worked 5 hours, one who worked 20, and one who
worked 40. Which of them could end up at a kink?
\end{exercise}

\begin{exercise}[$\star\star\star$]\label{exo:b1:constrained-optimisation:equivalence}
Let $f$ and $g$ be smooth, and let $(x^*, y^*)$ satisfy $g = b$ with
$g_x \neq 0$, $g_y \neq 0$ and $f_y \neq 0$.
\begin{enumerate}
  \item Show that if the gradients $\nabla f$ and $\nabla g$ are parallel at $(x^*, y^*)$, the \omterm{thm:b1:constrained-optimisation:lagrange}{Lagrange conditions} hold.
  \item Show that the \omterm{thm:b1:constrained-optimisation:lagrange}{Lagrange conditions} hold if and only if the level curve of $f$ and the constraint have a common tangent line.
  \item Give an example where the \omterm{thm:b1:constrained-optimisation:lagrange}{Lagrange conditions} hold at a \omterm{def:b1:constrained-optimisation:constrained-max}{constrained minimum}, and draw it.
\end{enumerate}
\end{exercise}

\begin{exercise}[$\star\star\star$]\label{exo:b1:constrained-optimisation:dual}
Minimise the cost $x + 2y$ of reaching the level curve $x^{1/2}y^{1/2} = u_0$.
\begin{enumerate}
  \item Solve with \omterm{met:b1:constrained-optimisation:lagrange}{Lagrange's method}; give the minimum cost $C(u_0)$ and the multiplier $\mu^*$.
  \item Show that $C'(u_0) = \mu^*$, and give its unit.
  \item Take $u_0 = \sqrt{18}$. Compare with \cref{ex:b1:constrained-optimisation:substitution}, and show that $\mu^* = 1/\lambda^*$.
\end{enumerate}
\end{exercise}

\begin{exercise}[$\star\star\star$]\label{exo:b1:constrained-optimisation:falsify}
Someone reads the opening chart as proof that ``when the wage rises, a worker
works less because the income effect beats the substitution effect.''
\begin{enumerate}
  \item What would the chart look like if the reading were false and every worker had a rising labour supply curve?
  \item Name two differences between rich and poor countries, other than the wage, that could produce the falling cloud.
  \item What kind of data, about the same people over time, could contradict the reading?
\end{enumerate}
\end{exercise}

\section{Problem: A Farm with One Field}\label{sec:b1:constrained-optimisation:problem}

\begin{problem}[{Weekend problem --- the shadow rent of a hectare}]\label{pb:b1:constrained-optimisation:farm}
A farm near Brenmoor (invented) has a field of 20 hectares to share between
potatoes and barley. On $x$ hectares of potatoes the season's profit is
$2400\sqrt x$ dalens, and on $y$ hectares of barley $1200\sqrt y$ dalens; the
square roots say that the farmer plants the best soil first. The farmer
maximises total profit $\Pi(x, y) = 2400\sqrt x + 1200\sqrt y$ subject to
$x + y = 20$, $x, y \ge 0$.

\medskip\noindent\textbf{Part I --- By substitution.}
\begin{enumerate}
  \item Write profit as a function of $x$ alone.
  \item Find the $x$ that makes its derivative zero, and the matching $y$.
  \item Check that it is a maximum, using the second derivative.
  \item Compute the maximum profit.
  \item Compare with planting the whole field in potatoes, and in barley.
  \item Compute the marginal profit of a hectare of potatoes and of barley at the optimum. Why must they be equal?
\end{enumerate}

\medskip\noindent\textbf{Part II --- By Lagrange's method.}
\begin{enumerate}[resume]
  \item Write the \omterm{def:b1:constrained-optimisation:lagrangian}{Lagrangian}.
  \item Write the three conditions.
  \item Divide the first two to obtain the \omterm{def:b1:constrained-optimisation:tangency}{tangency condition}, and solve it for $x/y$.
  \item Solve for $x$, $y$ and $\lambda$.
  \item Draw the constraint and the level curve of $\Pi$ through the optimum, and mark the optimum.
  \item Explain why \omterm{met:b1:constrained-optimisation:lagrange}{Lagrange's method} and substitution must give the same point here.
\end{enumerate}

\medskip\noindent\textbf{Part III --- The shadow price of land.}
\begin{enumerate}[resume]
  \item Solve the problem with $L$ hectares in place of 20, and show that the \omterm{def:b1:constrained-optimisation:value-function}{value function} is $V(L) = 6000\sqrt{L/5}$.
  \item Compute $V'(20)$ and compare it with $\lambda$.
  \item Give the unit of $\lambda$ and its meaning.
  \item A neighbour offers to rent one hectare to the farm for \dal{250} for the season. Should the farm, maximising profit, take it? And at \dal{350}?
  \item Compute the exact gain from a 21st hectare, $V(21) - V(20)$, and compare it with $\lambda$.
  \item Does the shadow rent rise or fall with the size of the field? Explain from $V(L)$.
\end{enumerate}

\medskip\noindent\textbf{Part IV --- A second constraint.}
Potatoes need 500 cubic metres of water a hectare in the season, barley
100. The farm has a water allowance of $W$ cubic metres.
\begin{enumerate}[resume]
  \item How much water does the plan of Part II use? Is the water constraint \omterm{def:b1:constrained-optimisation:binding}{binding} when $W = 9000$? What is its \omterm{def:b1:constrained-optimisation:shadow-price}{shadow price} then?
  \item Take $W = 6000$. Show that both constraints bind, and find $x$ and $y$ from the two constraints.
  \item Compute profit, and the profit lost because of the water allowance.
  \item Write the \omterm{def:b1:constrained-optimisation:lagrangian}{Lagrangian} with two multipliers, $\lambda$ for land and $\mu$ for water, and its conditions at the new plan.
  \item Solve for $\mu$ and $\lambda$.
  \item State the chapter's result: the shadow rent of one hectare when water is plentiful and when it is scarce, and what the farm would pay for one more cubic metre of water.
\end{enumerate}
\end{problem}
